% =====================================================================
% Relativistic Extension of Chapter III, §§24–28
% Stationary Convection with Rotation
% =====================================================================
% Agent 13 — Relativistic modifications of Chandrasekhar's treatment
% of thermal instability in a rotating fluid layer, focusing on
% stationary (non-oscillatory) convection.
%
% Conventions: see RELATIVISTIC_CONVENTIONS.md and rel_preamble.tex
% Metric signature: (-,+,+,+); c kept explicit; Israel–Stewart dissipation.
% =====================================================================

\section{Thermal instability in a rotating relativistic fluid: general considerations}
\label{sec:rel-3-24}

The classical arguments of \S\,24 carry over to the relativistic setting
once we replace the Newtonian equations of motion by the conservation law
$\nabla_\nu T^{\mu\nu} = 0$ for a dissipative fluid in a gravitational
field, augmented by the Coriolis coupling that arises in a frame
co-rotating with angular velocity $\boldsymbol{\Omega}$.

In the non-relativistic theory the Taylor--Proudman theorem forbids
steady motions with spatial variation along $\boldsymbol{\Omega}$ in an
inviscid fluid, so that only viscous fluids can convect.  The relativistic
analogue proceeds as follows.  In a stationary ($\partial_t = 0$),
linearised flow the spatial part of $\nabla_\nu T^{\mu\nu} = 0$ in the
co-rotating frame gives, at leading order,
\begin{equation}\label{eq:rel-3-TP}
  2\,\frac{\enthalpy}{c^2}\,\epsilon_{ijk}\,\Omega_j\,\delta u_k
  = -\partial_i \delta p + \rdensity\,g\,\alpha\,\Theta\,\lambda_i
    + \shearvisc\,\nabla^2 \delta u_i
    + \order{v^2/c^2}\;\text{corrections},
\end{equation}
where $\enthalpy = \edensity + p$ is the relativistic enthalpy density,
$\delta u_i$ is the spatial velocity perturbation, $\Theta$ the
temperature perturbation, $\alpha$ the thermal expansion coefficient,
and $\lambda_i = (0,0,1)$.

When $\shearvisc = 0$ and $\Theta = 0$ the equation reduces to the
relativistic Taylor--Proudman constraint: the Coriolis term on the left
(proportional to $\enthalpy/c^2$ rather than $\rho_0$) must be balanced
by the pressure gradient alone, forbidding velocity variation along
$\boldsymbol{\Omega}$.  The replacement $\rho_0 \to \enthalpy/c^2$ is the
characteristic relativistic modification; in a radiation-dominated fluid
where $p = \edensity/3$ and $\enthalpy = 4\edensity/3$, this enhancement
is significant.

\begin{relcorrection}
The relativistic Taylor--Proudman inertia is $\enthalpy/c^2$ rather than
$\rdensity$.  For matter with equation of state $p = \edensity/3$
(radiation-dominated), $\enthalpy/c^2 = 4\edensity/(3c^2)$, which
exceeds $\rdensity$ and thereby stiffens the rotational constraint.
The net effect is to \emph{increase} the stabilising influence of
rotation beyond the Newtonian prediction.
\end{relcorrection}

As in the classical case, the Taylor--Proudman constraint does not
forbid oscillatory motions: inertial (Coriolis) waves propagate with a
dispersion relation that receives relativistic corrections of order
$v^2/c^2$ (and additional Israel--Stewart dissipative corrections; see
below).  Consequently, overstability may again be the preferred route to
instability under certain conditions; this will be addressed in a
companion section on oscillatory convection.

% =====================================================================
\section{The relativistic perturbation equations with rotation}
\label{sec:rel-3-25}

\subsection*{Setup and governing equations}

We consider an infinite horizontal layer of relativistic fluid of
depth~$d$, confined between $z = 0$ and $z = d$, rotating uniformly
with angular velocity $\boldsymbol{\Omega} = \Omega\,\boldsymbol{e}_z$
(parallel to gravity $\mathbf{g} = -g\,\boldsymbol{e}_z$).  The background
is a static state with a steady adverse temperature gradient $\beta > 0$.

In the co-rotating frame the relativistic Euler equation for a
dissipative fluid, expanded to first order in perturbations and to
the relevant post-Newtonian order, reads (cf.\ the classical
equation~\eqref{eq:3-78})
\begin{equation}\label{eq:rel-3-78}
  \frac{\enthalpy}{c^2}\,\frac{\partial u_i}{\partial t}
  = -\frac{\partial(\delta p)}{\partial x_i}
    + \rdensity\,g\,\alpha\,\Theta\,\lambda_i
    + \shearvisc\,\nabla^2 u_i
    + 2\,\frac{\enthalpy}{c^2}\,\epsilon_{ijk}\,u_j\,\Omega_k
    + \frac{1}{\tauq}\,\partial_i q_{\mathrm{IS}},
\end{equation}
where the last term encodes the Israel--Stewart heat-flux correction
(see below).  The factor $\enthalpy/c^2$ multiplying the inertial and
Coriolis terms replaces the Newtonian $\rho_0$.

The heat equation, in the Israel--Stewart causal framework, becomes
a telegraph-type equation rather than the parabolic Fourier law:
\begin{equation}\label{eq:rel-3-79}
  \tauq\,\frac{\partial^2\Theta}{\partial t^2}
  + \frac{\partial\Theta}{\partial t}
  = \beta\,w + \kappa\,\nabla^2\Theta
  + \tauq\,\kappa\,\frac{\partial}{\partial t}\nabla^2\Theta,
\end{equation}
where $\tauq$ is the heat-flux relaxation time.  In the limit
$\tauq \to 0$ this reduces to the classical equation~\eqref{eq:3-79}.

The continuity equation retains its classical form at leading order:
\begin{equation}\label{eq:rel-3-80}
  \frac{\partial u_j}{\partial x_j} = 0.
\end{equation}

Taking the curl of~\eqref{eq:rel-3-78} (exactly as in the classical
derivation, \S\,25) and projecting onto $\boldsymbol{\lambda}$, we obtain
the relativistic analogues of the vorticity and vertical-velocity
equations:
\begin{align}
  \frac{\enthalpy}{c^2}\,\frac{\partial\zeta}{\partial t}
  &= \shearvisc\,\nabla^2\zeta
    + 2\,\frac{\enthalpy}{c^2}\,\Omega\,\frac{\partial w}{\partial z},
  \label{eq:rel-3-87}\\[6pt]
  \frac{\enthalpy}{c^2}\,\frac{\partial}{\partial t}\nabla^2 w
  &= \rdensity\,g\,\alpha\!\left(
       \frac{\partial^2\Theta}{\partial x^2}
       + \frac{\partial^2\Theta}{\partial y^2}\right)
    + \shearvisc\,\nabla^4 w
    - 2\,\frac{\enthalpy}{c^2}\,\Omega\,\frac{\partial\zeta}{\partial z}.
  \label{eq:rel-3-88}
\end{align}
Here $w$ and $\zeta$ are the $z$-components of velocity and vorticity.

\subsection*{(a) Normal mode analysis with Israel--Stewart dissipation}

We decompose perturbations into normal modes of the form
$\propto \exp(pt + i\mathbf{k}_\perp\cdot\mathbf{x}_\perp)$ with
horizontal wavenumber $k$ and growth rate $p$.  Using
the non-dimensionalisation $a = kd$, $\sigma = pd^2/\nu_{\mathrm{eff}}$
where $\nu_{\mathrm{eff}} = \shearvisc\,c^2/\enthalpy$ is the
relativistic kinematic viscosity, and $\mathscr{P}_{\mathrm{rel}} =
\nu_{\mathrm{eff}}/\kappa$, we obtain the relativistic counterparts of
equations~\eqref{eq:3-93}--\eqref{eq:3-95}:
\begin{align}
  \bigl[(1 + \tauq\, p)\,(D^2 - a^2)
     - \mathscr{P}_{\mathrm{rel}}\,\sigma\bigr]\,\Theta
  &= -\frac{\beta d^2}{\kappa}\,W,
  \label{eq:rel-3-93}\\[6pt]
  (D^2 - a^2 - \sigma)\,Z
  &= -\frac{2\Omega}{\nu_{\mathrm{eff}}}\,d\;DW,
  \label{eq:rel-3-94}\\[6pt]
  (D^2 - a^2)(D^2 - a^2 - \sigma)\,W
  - \frac{2\Omega}{\nu_{\mathrm{eff}}}\,d^3\;DZ
  &= \frac{\rdensity\,g\,\alpha\,d^2}{\shearvisc}\;a^2\,\Theta.
  \label{eq:rel-3-95}
\end{align}

The factor $(1 + \tauq\,p)$ in the temperature
equation~\eqref{eq:rel-3-93} is the hallmark of Israel--Stewart
causal heat conduction.  It converts the parabolic heat equation into
a damped wave equation, bounding the thermal signal speed by
\begin{equation}\label{eq:rel-heat-speed}
  v_{\mathrm{heat}}^2
  = \frac{\kappa}{\tauq}
  \leq c^2.
\end{equation}
The inequality $\tauq \geq \kappa/c^2$ is required for causality.

\begin{causalitycheck}
The Israel--Stewart relaxation time $\tauq$ ensures that the heat mode
propagates at speed $v_{\mathrm{heat}} = \sqrt{\kappa/\tauq} \leq c$.
In the normal-mode equations, the factor $(1+\tauq\,p)$ regularises
the spectrum: for $|p| \to \infty$ the temperature equation becomes
$\tauq\,p\,(D^2 - a^2)\Theta \approx -(\beta d^2/\kappa)\,W$, giving
a finite propagation speed rather than the instantaneous diffusion of
the Fourier law.  All perturbation modes therefore respect causality.
\end{causalitycheck}


% =====================================================================
\section{The case of stationary convection: a relativistic variational principle}
\label{sec:rel-3-26}

When instability sets in as stationary convection the marginal state has
$\sigma = 0$ (and hence $p = 0$), so the Israel--Stewart time-derivative
terms drop out.  The marginal equations reduce to
\begin{align}
  (D^2 - a^2)\,\Theta
  &= -\frac{\beta d^2}{\kappa}\,W,
  \label{eq:rel-3-96}\\[4pt]
  (D^2 - a^2)\,Z
  &= -\frac{2\Omega}{\nu_{\mathrm{eff}}}\,d\;DW,
  \label{eq:rel-3-97}\\[4pt]
  (D^2 - a^2)^2 W
  - \frac{2\Omega}{\nu_{\mathrm{eff}}}\,d^3\;DZ
  &= \frac{\rdensity\,g\,\alpha\,d^2}{\shearvisc}\;a^2\,\Theta.
  \label{eq:rel-3-98}
\end{align}

Eliminating $Z$ and $\Theta$ yields the single equation for~$W$:
\begin{equation}\label{eq:rel-3-99}
  (D^2 - a^2)^3\,W + \Tarel\;D^2 W = -\Rarel\;a^2\,W,
\end{equation}
where the \emph{relativistic Rayleigh and Taylor numbers} are defined by
\begin{equation}\label{eq:rel-Ra-Ta-def}
  \Rarel
  = \frac{\rdensity\,g\,\alpha\,\beta\,d^4}{\shearvisc\,\kappa}
    \;\relcorr{v^2/c^2},
  \qquad
  \Tarel
  = \frac{4\,\Omega^2\,d^4}{\nu_{\mathrm{eff}}^2}
  = \frac{4\,\Omega^2\,d^4\,\enthalpy^2}{\shearvisc^2\,c^4}.
\end{equation}

\begin{relcorrection}
The relativistic Taylor number $\Tarel$ differs from the classical
$\mathscr{T} = 4\Omega^2 d^4/\nu^2$ by the replacement
$\nu \to \nu_{\mathrm{eff}} = \shearvisc\,c^2/\enthalpy$.  Since
$\enthalpy/c^2 \geq \rdensity$, we have $\nu_{\mathrm{eff}} \leq \nu$,
and therefore $\Tarel \geq \mathscr{T}$.  \emph{Relativistic inertia
enhances the effective Taylor number.}  For a radiation-dominated gas
($p = \edensity/3$) the enhancement factor is
$(4\edensity/3\rdensity c^2)^2$, which can be substantial in
astrophysical contexts (e.g.\ the cores of proto-neutron stars).
\end{relcorrection}

The relativistic Rayleigh number receives corrections from two sources:
(i) the replacement of $\rho_0$ by $\enthalpy/c^2$ in the viscous
diffusion time, and (ii) modifications to the buoyancy arising from the
relativistic equation of state $\delta\rho/\rho_0 =
-\alpha\,\Theta\,[1 + \order{p/\edensity}]$.  These corrections are
absorbed into the $\relcorr{v^2/c^2}$ factor in~\eqref{eq:rel-Ra-Ta-def}.

\subsection*{(a) A variational principle}

The self-adjoint structure of equations~\eqref{eq:rel-3-97}
and~\eqref{eq:rel-3-99} is preserved under the relativistic
modifications (the replacement $\nu \to \nu_{\mathrm{eff}}$ does not
break the symmetry).  Defining
\begin{equation}\label{eq:rel-3-104}
  F = (D^2 - a^2)^2 W
    - \frac{2\Omega}{\nu_{\mathrm{eff}}}\,d^3\;DZ,
\end{equation}
the variational expression for the relativistic Rayleigh number is
\begin{equation}\label{eq:rel-3-114}
  \Rarel
  = \frac{\displaystyle\int_0^1 \bigl[(DF)^2 + a^2 F^2\bigr]\,dz}
         {a^2\displaystyle\int_0^1
           \bigl\{[(D^2-a^2)W]^2
             + d^2[(DZ)^2 + a^2 Z^2]\bigr\}\,dz}
  \equiv \frac{I_1}{a^2\,I_2},
\end{equation}
where now all instances of $\nu$ in the underlying equations are
replaced by $\nu_{\mathrm{eff}}$.

The proof of stationarity ($\delta\Rarel = 0$ iff the Euler--Lagrange
equation holds) follows line-by-line from the classical derivation in
\S\,26\,(a), since only algebraic and integration-by-parts identities
are used; the specific value of the viscosity parameter is immaterial.
Consequently, the lowest eigenvalue of $\Rarel$ is a true minimum of the
functional~\eqref{eq:rel-3-114}, and the same variational procedure
described in \S\,26 can be applied in the relativistic case.


% =====================================================================
\section{Solutions for stationary convection in a rotating relativistic fluid}
\label{sec:rel-3-27}

\subsection*{(a) Two free boundaries: relativistic corrections to critical $\Rarel(\Tarel)$}

With both boundaries free the boundary conditions remain the same as in
the classical case~\eqref{eq:3-124}, and the eigenfunctions are
$W = A\sin n\pi z$.  The characteristic equation becomes
\begin{equation}\label{eq:rel-3-127}
  \Rarel
  = \frac{1}{a^2}\bigl[(n^2\pi^2 + a^2)^3 + n^2\pi^2\,\Tarel\bigr].
\end{equation}

For the lowest mode ($n=1$), minimising over $a^2 = \pi^2 x$ gives
\begin{equation}\label{eq:rel-3-131}
  2x^3 + 3x^2 = 1 + \Tarel/\pi^4.
\end{equation}

This is identical in form to the classical equation~\eqref{eq:3-131},
but with $\Tarel$ replacing $\mathscr{T}$.  Since $\Tarel \geq
\mathscr{T}$, the critical Rayleigh number is \emph{higher} than the
Newtonian prediction at the same physical rotation rate: relativistic
inertia strengthens the rotational stabilisation.

The critical values are:
\begin{equation}\label{eq:rel-3-128}
  \Rarel^{(c)}
  = \pi^4\left[(1+x_{\min})^3 + \frac{\Tarel}{\pi^4}\right]
    \frac{1}{x_{\min}},
\end{equation}
where $x_{\min}$ is the appropriate root of~\eqref{eq:rel-3-131}.

To translate back to the physical critical temperature gradient:
\begin{equation}\label{eq:rel-beta-c}
  g\,\alpha\,\beta_c
  = \frac{\Rarel^{(c)}\;\shearvisc\,\kappa}{\rdensity\,d^4}
    \;\relcorr{v^2/c^2}^{-1}.
\end{equation}

\subsection*{(b) Two rigid boundaries: modified Taylor number $\Tarel$}

For rigid boundaries the variational method of \S\,26 applies without
structural change.  We expand $F = \sum_m A_m \cos[(2m+1)\pi z]$ and
solve the eighth-order system exactly as in the classical \S\,27\,(b),
with $\nu \to \nu_{\mathrm{eff}}$ throughout.

The characteristic equation is
\begin{equation}\label{eq:rel-3-159}
  \left\|\frac{1}{2}\,c_{m+1}
    \left[\frac{1}{\Rarel\,a^2} - \gamma_{2m+1}\right]
    \delta_{mn} - (n|m)\right\| = 0,
\end{equation}
where $c_{m+1}$, $\gamma_{2m+1}$, and $(n|m)$ are defined exactly as
in equations~\eqref{eq:3-139}, \eqref{eq:3-147}, and~\eqref{eq:3-165},
but the cubic~\eqref{eq:3-148} now involves $\Tarel$:
\begin{equation}\label{eq:rel-3-148}
  (q^2 - a^2)^3 + \Tarel\;q^2 = 0.
\end{equation}

The numerical tables (Tables~VII--IX in the original) can be reused
directly by reading off values at $\mathscr{T} = \Tarel$ rather than at
the Newtonian $\mathscr{T}$.  Since $\Tarel > \mathscr{T}$ for any
relativistic fluid, \emph{the critical Rayleigh number for rigid
boundaries is shifted upward compared to the Newtonian prediction at the
same physical rotation rate}.

\subsection*{(c) One rigid and one free boundary}

The reduction of this case to the odd solutions of the two-rigid-boundary
problem (\S\,27\,(c) in the original) carries over unchanged: an odd
solution for $W$ on $[-\tfrac12,+\tfrac12]$ satisfies free-boundary
conditions at $z = 0$ and rigid conditions at $z = \pm\tfrac12$.  The
Rayleigh and Taylor numbers for a cell of depth~$d$ are $\tfrac{1}{16}$
of those for a cell of depth~$\tfrac12 d$ in the two-rigid problem.
All relativistic modifications enter solely through $\Tarel$.

\subsection*{(d) The $\mathscr{T}^{1/3}$- and $\mathscr{T}^{2/3}$-laws in the relativistic regime}

We now examine whether the classical asymptotic scaling laws persist
when $\Tarel \to \infty$.

Starting from the full equation~\eqref{eq:rel-3-99} and retaining only
the leading terms in $\Tarel$ and the highest power of~$a$
(following the argument of classical \S\,27\,(d)), we obtain
\begin{equation}\label{eq:rel-3-179}
  \Tarel\;D^2 W = -(\Rarel\,a^2 - a^6)\,W.
\end{equation}

This is identical in structure to equation~\eqref{eq:3-179}; the
eigenvalue is $W = A\sin n\pi z$ with
\begin{equation}\label{eq:rel-3-181}
  \Rarel = \frac{1}{a^2}(a^6 + \pi^2\,\Tarel).
\end{equation}

Minimising over $a$:
\begin{equation}\label{eq:rel-3-183}
  a_{\min} = (\tfrac{1}{2}\pi^2\,\Tarel)^{1/6},
  \qquad
  \Rarel^{(c)} = 3(\tfrac{1}{2}\pi^2\,\Tarel)^{2/3}
               = 8.6956\;\Tarel^{2/3}.
\end{equation}

\begin{relcorrection}
The $\Tarel^{1/3}$- and $\Tarel^{2/3}$-laws persist in the
relativistic theory.  The power-law exponents $1/6$ and $2/3$ are
\emph{universal} --- they follow from dimensional analysis of the
reduced second-order equation and are insensitive to the specific value
of the viscosity.  However, the physical content changes:
\[
  \Tarel^{2/3}
  = \left(\frac{4\Omega^2 d^4}{\nu_{\mathrm{eff}}^2}\right)^{\!2/3}
  = \left(\frac{4\Omega^2 d^4\,\enthalpy^2}
              {\shearvisc^2 c^4}\right)^{\!2/3}.
\]
In terms of the \emph{physical} rotation rate, the critical temperature
gradient at large $\Omega$ is
\begin{equation}\label{eq:rel-3-134}
  g\,\alpha\,\beta_c
  \;\to\; 21.911\,\left(\frac{\Omega}{d}\right)^{\!4/3}
          \kappa^{-1}
          \left(\frac{\enthalpy}{\rdensity\,c^2}\right)^{\!4/3}
  \quad (\Omega \to \infty),
\end{equation}
which exceeds the classical result~\eqref{eq:3-134} by the factor
$(\enthalpy/\rdensity c^2)^{4/3} \geq 1$.  For non-relativistic matter
($\enthalpy \approx \rdensity c^2$) this reduces to the classical
formula.
\end{relcorrection}

The origin of the scaling laws is the same as in the classical theory:
for $\Tarel \to \infty$ the effective order of the governing ODE drops
from six to two, and only two boundary conditions ($W = 0$ at both
surfaces) can be satisfied.  The relativistic modifications do not alter
the order reduction; they only modify the prefactors through
$\nu_{\mathrm{eff}}$.


% =====================================================================
\section{Cell patterns with rotation: relativistic modifications}
\label{sec:rel-3-28}

\subsection*{Horizontal motions}

In terms of the solutions for $w$ and $\zeta$, the horizontal velocity
components retain the classical form~\eqref{eq:3-185}:
\begin{equation}\label{eq:rel-3-185}
  u = \frac{1}{a^2}\!\left(
    \frac{\partial^2 w}{\partial x\,\partial z}
    + \frac{\partial\zeta}{\partial y}\right), \quad
  v = \frac{1}{a^2}\!\left(
    \frac{\partial^2 w}{\partial y\,\partial z}
    - \frac{\partial\zeta}{\partial x}\right).
\end{equation}

The relationship between $Z$ and $DW$ is now
\begin{equation}\label{eq:rel-3-188}
  (D^2 - a^2)\,Z
  = -\frac{2\Omega\,d}{\nu_{\mathrm{eff}}}\;DW,
\end{equation}
so the coupling coefficient involves $\nu_{\mathrm{eff}}$ rather
than~$\nu$.  Since $\nu_{\mathrm{eff}} \leq \nu$, the Coriolis coupling
to horizontal vorticity is \emph{stronger} in the relativistic case.

For two free boundaries the lowest-mode solutions are
\begin{equation}\label{eq:rel-3-192}
  W = W_0\sin\pi z, \qquad
  Z = W_0\,\frac{2\Omega d}{\nu_{\mathrm{eff}}}
      \,\frac{\pi}{\pi^2+a^2}\,\cos\pi z.
\end{equation}

The mean horizontal kinetic energy is
\begin{equation}\label{eq:rel-3-193}
  \langle u^2 + v^2 \rangle
  = \tfrac14 W_0^2\;
    \frac{(\pi^2+a^2)^2 + \Tarel}{a^2(\pi^2+a^2)^2}\;\cos^2\pi z.
\end{equation}

At marginal stability, using~\eqref{eq:rel-3-131} to
eliminate~$\Tarel$, we recover
\begin{equation}\label{eq:rel-3-195}
  \langle u^2 + v^2 \rangle
  = \tfrac14 W_0^2\,\cos^2\pi z,
\end{equation}
which is \emph{independent} of $\Tarel$, exactly as in the classical
theory.  The remarkable invariance of the ratio of horizontal to
vertical kinetic energies at the onset of stationary convection
thus persists in the relativistic regime.

\subsection*{Preferred cell shapes}

The cell patterns (rolls, squares, hexagons) analysed by Veronis in the
classical theory are determined by the horizontal planform of the
eigenfunctions, which is independent of the vertical structure and hence
independent of whether $\nu$ or $\nu_{\mathrm{eff}}$ appears in the
equations.  The explicit formulae for the horizontal velocities
(equations~\eqref{eq:3-196}, \eqref{eq:3-204}, \eqref{eq:3-207}) carry
over with the single substitution
\begin{equation}\label{eq:rel-T-sub}
  \sqrt{\mathscr{T}} \;\longrightarrow\; \sqrt{\Tarel}
  \quad\text{wherever it appears in the coupling to }\zeta.
\end{equation}

\paragraph{Rolls.}
The ratio $v/u$ becomes (cf.~\eqref{eq:3-200})
\begin{equation}\label{eq:rel-3-200}
  \frac{v}{u}
  = -\frac{\sqrt{\Tarel}}{\pi^2 + a^2},
\end{equation}
and the streamline wave number measured in the plane containing
the streamlines remains
\begin{equation}\label{eq:rel-3-203}
  a_s^2 = \tfrac14\pi^2,
\end{equation}
independent of rotation and of relativistic corrections.  Veronis's
result that the wavelength of a roll measured in the streamline plane is
invariant under rotation extends to the relativistic case.

\paragraph{Square and hexagonal cells.}
The spiral patterns caused by Coriolis acceleration are tighter in the
relativistic case because $\sqrt{\Tarel} > \sqrt{\mathscr{T}}$ at the
same physical $\Omega$.  The qualitative features --- fluid spiralling
clockwise (for counter-clockwise $\Omega$) as it rises, reversing at
mid-plane, and spiralling counter-clockwise as it descends --- are
unchanged.

\paragraph{Limiting behaviour as $\Tarel \to \infty$.}
The dominant horizontal motions become
\begin{equation}\label{eq:rel-3-208}
  u \to -\frac{\pi\sqrt{\Tarel}}{a^2(\pi^2+a^2)}\,a_y\,W_0
        \cos a_x x\,\sin a_y y\,\cos\pi z,
  \quad
  v \to +\frac{\pi\sqrt{\Tarel}}{a^2(\pi^2+a^2)}\,a_x\,W_0
        \sin a_x x\,\cos a_y y\,\cos\pi z,
\end{equation}
so the streamlines are $w = \mathrm{const}$ on horizontal planes,
exactly as in the non-relativistic limit~\eqref{eq:3-210}.  The small
radial component proportional to $\nabla_\perp w$ is responsible for
convective transport when $\Tarel$ is large.

\begin{relcorrection}
Quantitative differences from the classical cell patterns are:
\begin{enumerate}
  \item The spirals are wound more tightly by a factor
        $\sqrt{\Tarel/\mathscr{T}} = \enthalpy/(\rdensity c^2)$
        at the same physical rotation rate.
  \item The transition to the asymptotic regime of nearly closed
        streamlines occurs at a \emph{lower} physical rotation rate
        $\Omega$ because $\Tarel$ grows faster with $\enthalpy$.
  \item For non-relativistic matter the corrections are
        $\order{p/\edensity}$, typically negligible.  For
        radiation-dominated or degenerate matter the corrections
        can be of order unity.
\end{enumerate}
\end{relcorrection}


% =====================================================================
\subsection*{Causality verification for rotating stationary convection}

\begin{causalitycheck}
We verify that all perturbation modes in the rotating system respect
causality ($v_{\mathrm{signal}} \leq c$).

\textbf{1. Heat propagation.}
The Israel--Stewart heat equation~\eqref{eq:rel-3-79} has characteristic
speed $v_{\mathrm{heat}} = \sqrt{\kappa/\tauq}$.  The causal constraint
$\tauq \geq \kappa/c^2$ ensures $v_{\mathrm{heat}} \leq c$.
At marginal stability ($p = 0$) the time-derivative terms vanish, so
causality is automatically satisfied for the stationary mode.

\textbf{2. Viscous (shear) modes.}
In the Israel--Stewart framework the shear stress obeys
$\taupi\,\dot\pi^{\langle\mu\nu\rangle} + \pi^{\mu\nu} =
2\shearvisc\,\sigma^{\mu\nu}$, giving a shear signal speed
$v_{\mathrm{shear}} = \sqrt{\shearvisc/(\enthalpy\,\taupi/c^2)}$.
The constraint $\taupi \geq \shearvisc\,c^2/\enthalpy\cdot c^{-2}
= \shearvisc/\enthalpy$ ensures $v_{\mathrm{shear}} \leq c$.
At marginal stability these modes are not excited.

\textbf{3. Inertial (Coriolis) waves.}
The Coriolis coupling produces inertial waves with dispersion relation
$\omega^2 = 4\Omega^2\cos^2\theta\;[1 + \order{v^2/c^2}]$, where
$\theta$ is the angle between $\mathbf{k}$ and $\boldsymbol{\Omega}$.
Since $\Omega \ll c/d$ in any physically realisable rotating fluid
layer, the group velocity $v_g \sim \Omega/k \ll c$.  Causality is not
threatened.

\textbf{4. Combined modes.}
The full dispersion relation obtained from
equations~\eqref{eq:rel-3-93}--\eqref{eq:rel-3-95} is a polynomial in
$\sigma$ with finite-order terms only (the Israel--Stewart relaxation
times ensure that the highest power of $\sigma$ has a positive
coefficient).  By the standard argument for symmetric hyperbolic systems,
all characteristic speeds are bounded by the maximum of the individual
signal speeds, which are all $\leq c$.

\textbf{Conclusion:} All perturbation modes in the rotating
relativistic convection problem are causal, both at marginal stability
and in the general time-dependent regime.
\end{causalitycheck}


% =====================================================================
\subsection*{Summary of relativistic modifications to \S\S\,24--28}

\begin{enumerate}
  \item The Taylor--Proudman theorem carries over with the inertial mass
        density $\rdensity \to \enthalpy/c^2$.
  \item The perturbation equations (\S\,25) acquire the Israel--Stewart
        causal heat equation and the replacement
        $\nu \to \nu_{\mathrm{eff}} = \shearvisc\,c^2/\enthalpy$
        in all Coriolis and viscous terms.
  \item The variational principle (\S\,26) retains its form; the
        self-adjointness is not broken.
  \item For stationary convection (\S\,27), the eigenvalue problem is
        structurally identical to the classical one but with
        $\mathscr{T} \to \Tarel \geq \mathscr{T}$.  All classical
        tables and figures can be reused at the shifted Taylor number.
  \item The $\Tarel^{1/3}$ and $\Tarel^{2/3}$ asymptotic laws persist
        with modified prefactors containing $\enthalpy/\rdensity c^2$.
  \item Cell patterns (\S\,28) are qualitatively unchanged; spirals are
        wound more tightly by the factor $\sqrt{\Tarel/\mathscr{T}}$.
  \item All modes are causal when Israel--Stewart relaxation times
        satisfy the standard lower bounds.
  \item In the non-relativistic limit
        $\enthalpy \to \rdensity c^2$, $\tauq \to 0$, all results
        reduce to those of Chandrasekhar's original \S\S\,24--28.
\end{enumerate}
